% !TeX root = ../main.tex
\section{Next steps}
\label{sec:next}

\subsection{Order of work}

The results point at the dry signal first: the search is close to optimal when
it has the right one.

\begin{nextsteps}
\begin{enumerate}
\item \textbf{A better dry signal} \wip: the sampler of real DI notes; later a
  learned DI renderer (in the manner of GuitarFlow~\cite{loth2025guitarflow},
  trained on redistributable DI), an effect remover trained on the app's own
  renders~\cite{lee2024distortion,yeh2026egvae}, and the player's own DI through a
  pre-chain tap in practice mode.
\item \textbf{Separated stems} in the benchmark \wip, then the transcription
  pipeline's own notes instead of labels.
\item \textbf{Time-based effects from notes and tempo}: delay \wip; reverb decay
  from rests; modulation rate from sustained notes \planned.
\item \textbf{A learned initial guess} trained on separated renders of the
  template chain, returning several candidates \planned.
\item \textbf{Listening tests} (MUSHRA via webMUSHRA) to check that the distance
  agrees with ears \planned.
\item \textbf{The UI}: save as a preset tagged with the song and track \wip;
  compare the preset with the stem played in sync by the tab player \planned.
\item \textbf{Captures}: the user's NAM files as amp candidates \planned.
\end{enumerate}
\end{nextsteps}

\subsection{The full evaluation protocol}

The benchmark is meant to grow into four sets:
\begin{enumerate}
\item \textbf{In-chain synthetic}: held-out DI (Guitar-TECHS P3, GuitarSet
  players 03--05) through random template chains, isolated and over backing tracks
  after separation; parameter error and audio distances. The difference between
  isolated and separated inputs measures the cost of separation.
\item \textbf{Equipment outside the chain}: EGDB and EGDB-PG renders, Guitar-TECHS
  real amplifiers, and for research only GOAT and ToneTwist devices; the true DI
  is rendered through the estimate and compared with the real output.
\item \textbf{Real productions}: MoisesDB songs with isolated guitar tracks;
  estimates from the separated stem against estimates from the isolated track,
  each scored against the isolated track.
\item \textbf{Listening}: MUSHRA on sets 1 and 2, where the same performance
  exists on both sides; tone-similarity scores on set 3; correlation of the
  objective metrics with the ratings.
\end{enumerate}
Baselines throughout: the current metadata-based suggestion, the nearest factory
preset, that preset plus the match, the full pipeline, and the search with the
true DI as upper bound.

\subsection{Conclusion}

Each part of guitar tone estimation has working methods under easier conditions:
effects are recognised on synthetic isolated guitar, amplifiers are captured from
paired audio and, less accurately, from unpaired audio, tone-transfer models
reproduce whole chains as audio, and inference-time search drives arbitrary,
non-differentiable chains. What was missing is the combination the commercial
case needs: an editable chain, estimated from a separated stem without the dry
signal. The app now has a first version of that combination, using its own
renderer as the forward model, a gradient-free search, a final linear match,
and one idea the literature had not used: a dry signal synthesised from the
transcribed notes. The first benchmark shows that the search is not the weak
part; the stand-in dry signal is. Improving it, and measuring on separated stems
and with listeners, are the next steps.
